% GENERATED FILE. Edit canonical lessons, then run npm run generate:books.
% canonical-source-hash: fnv1a64:f32e5d0bdf1cc118
% canonical-lessons: TE-C285-bedirincu, TE-C285-heccarincu, TE-C285-gamanincu, TE-C285-tarumu, TE-C285-karugu

\chapter{To threaten, To warn, To take note of, To chase, To melt}
\label{ch:te-to-threaten-to-warn-to-take-note-of-to-chase-to-melt}
\hlchaptermodality{285}

\begin{chapteropening}
\textbf{By the end of this chapter:} \emph{I can say to threaten, to warn, to take note of, to chase and to melt.}

Complete the last lesson of chapter 285: I can say to threaten, to warn, to take note of, to chase and to melt.
\end{chapteropening}

\section[\texorpdfstring{\te{బెదిరించు} (bediriñcu)}{బెదిరించు (bediriñcu)}]{\te{బెదిరించు} (bediriñcu) — to threaten}

\label{lesson:TE-C285-bedirincu}



\begin{quote}
Before the new one: say the Telugu for to sweep, then the Telugu for to knead.
\end{quote}

\subsection*{You'll want to know: \te{బెదిరించు}}
\textbf{\te{బెదిరించు}} — \emph{bediriñcu} — \textquotedblleft{}to threaten\textquotedblright{}.

\begin{grammarlens}[title={What you've built}]
One more word for this chapter.
\end{grammarlens}

\subsection*{Your turn}
\begin{itemize}
\raggedright
  \item \emph{Say it:} \emph{bediriñcu}
  \item \emph{Say it:} \emph{bediriñcu}, once more
  \item \emph{Say it:} say \emph{bediriñcu}
  \item \emph{Recall:} say the Telugu for to sweep, then the Telugu for to knead, then say \emph{bediriñcu} again
\end{itemize}

\begin{tcolorbox}[breakable,colback=teal!4,colframe=teal!35!black,title={Before you move on}]
What does \te{బెదిరించు} mean? (\textquotedblleft{}To threaten\textquotedblright{}.) Say it once more.
\end{tcolorbox}
\section[\texorpdfstring{\te{హెచ్చరించు} (heccariñcu)}{హెచ్చరించు (heccariñcu)}]{\te{హెచ్చరించు} (heccariñcu) — to warn}

\label{lesson:TE-C285-heccarincu}



\begin{quote}
Before the new one: say the Telugu for to knead, then the Telugu for to threaten.
\end{quote}

\subsection*{You'll want to know: \te{హెచ్చరించు}}
\textbf{\te{హెచ్చరించు}} — \emph{heccariñcu} — \textquotedblleft{}to warn\textquotedblright{}.

\begin{grammarlens}[title={What you've built}]
One more word for this chapter.
\end{grammarlens}

\subsection*{Your turn}
\begin{itemize}
\raggedright
  \item \emph{Say it:} \emph{heccariñcu}
  \item \emph{Say it:} \emph{heccariñcu}, once more
  \item \emph{Say it:} say \emph{heccariñcu}
  \item \emph{Recall:} say the Telugu for to knead, then the Telugu for to threaten, then say \emph{heccariñcu} again
\end{itemize}

\begin{tcolorbox}[breakable,colback=teal!4,colframe=teal!35!black,title={Before you move on}]
What does \te{హెచ్చరించు} mean? (\textquotedblleft{}To warn\textquotedblright{}.) Say it once more.
\end{tcolorbox}
\section[\texorpdfstring{\te{గమనించు} (gamaniñcu)}{గమనించు (gamaniñcu)}]{\te{గమనించు} (gamaniñcu) — to take note of, to notice}

\label{lesson:TE-C285-gamanincu}



\begin{quote}
Before the new one: say the Telugu for to threaten, then the Telugu for to warn.
\end{quote}

\subsection*{You'll want to know: \te{గమనించు}}
\textbf{\te{గమనించు}} — \emph{gamaniñcu} — \textquotedblleft{}to take note of, to notice\textquotedblright{}.

\begin{grammarlens}[title={What you've built}]
One more word for this chapter.
\end{grammarlens}

\subsection*{Your turn}
\begin{itemize}
\raggedright
  \item \emph{Say it:} \emph{gamaniñcu}
  \item \emph{Say it:} \emph{gamaniñcu}, once more
  \item \emph{Say it:} say \emph{gamaniñcu}
  \item \emph{Recall:} say the Telugu for to threaten, then the Telugu for to warn, then say \emph{gamaniñcu} again
\end{itemize}

\begin{tcolorbox}[breakable,colback=teal!4,colframe=teal!35!black,title={Before you move on}]
What does \te{గమనించు} mean? (\textquotedblleft{}To take note of, to notice\textquotedblright{}.) Say it once more.
\end{tcolorbox}
\section[\texorpdfstring{\te{తరుము} (tarumu)}{తరుము (tarumu)}]{\te{తరుము} (tarumu) — to chase}

\label{lesson:TE-C285-tarumu}



\begin{quote}
Before the new one: say the Telugu for to warn, then the Telugu for to take note of.
\end{quote}

\subsection*{You'll want to know: \te{తరుము}}
\textbf{\te{తరుము}} — \emph{tarumu} — \textquotedblleft{}to chase\textquotedblright{}.

\begin{grammarlens}[title={What you've built}]
One more word for this chapter.
\end{grammarlens}

\subsection*{Your turn}
\begin{itemize}
\raggedright
  \item \emph{Say it:} \emph{tarumu}
  \item \emph{Say it:} \emph{tarumu}, once more
  \item \emph{Say it:} say \emph{tarumu}
  \item \emph{Recall:} say the Telugu for to warn, then the Telugu for to take note of, then say \emph{tarumu} again
\end{itemize}

\begin{tcolorbox}[breakable,colback=teal!4,colframe=teal!35!black,title={Before you move on}]
What does \te{తరుము} mean? (\textquotedblleft{}To chase\textquotedblright{}.) Say it once more.
\end{tcolorbox}
\section[\texorpdfstring{\te{కరుగు} (karugu)}{కరుగు (karugu)}]{\te{కరుగు} (karugu) — to melt}

\label{lesson:TE-C285-karugu}



\begin{quote}
Before the new one: say the Telugu for to take note of, then the Telugu for to chase.
\end{quote}

\subsection*{You'll want to know: \te{కరుగు}}
\textbf{\te{కరుగు}} — \emph{karugu} — \textquotedblleft{}to melt\textquotedblright{}.

\begin{grammarlens}[title={What you've built}]
One more word for this chapter.
\end{grammarlens}

\subsection*{Your turn}
\begin{itemize}
\raggedright
  \item \emph{Say it:} \emph{karugu}
  \item \emph{Say it:} \emph{karugu}, once more
  \item \emph{Say it:} say \emph{karugu}
  \item \emph{Recall:} say the Telugu for to take note of, then the Telugu for to chase, then say \emph{karugu} again
\end{itemize}

\begin{tcolorbox}[breakable,colback=teal!4,colframe=teal!35!black,title={Before you move on}]
What does \te{కరుగు} mean? (\textquotedblleft{}To melt\textquotedblright{}.) Say it once more.
\end{tcolorbox}
